% paper4_discussion.tex -- Discussion and Conclusion of Paper 4, input by
% papers/paper4_tillage.tex.
\section{Discussion}

Both reduced tillage and residue retention increased topsoil organic carbon,
and their combination gave the largest gain at every site. Residues supply the
carbon input; reduced disturbance slows the breakdown of aggregates that
protect organic matter from decomposition. Losses under conventional tillage
with residues removed show that, without either, cultivated soils continue to
lose carbon, consistent with the large carbon deficit of many agricultural
soils \parencite{lal2004}. Two cautions apply to these gains. They were
measured in the topsoil only, and no-till tends to concentrate carbon near the
surface, so part of the apparent gain may be a redistribution rather than a
net increase in the whole profile.

The yield results differed from the carbon results and from site to site. At
the temperate site, no-till lowered yield, particularly when residues were
removed; at the semi-arid site, no-till with residues retained gave the
highest yield; at the humid tropical site, neither factor mattered. This
pattern mirrors the conclusion of a global meta-analysis, in which no-till
reduced yields on average but performed best in dry climates and when combined
with residue retention \parencite{pittelkow2016}. Under dry conditions the
residue mulch conserves soil water, which can outweigh the yield penalty of
reduced tillage.

The results also show why conservation agriculture cannot be judged as a
single package. In many smallholder systems, residues are fed to livestock and
cannot simply be left on the field, and without residues the yield penalty of
no-till may be largest \parencite{giller2009}. The factorial design used here
separates the two components and shows that, at two of the three sites, the
benefit of no-till depended on the residues being retained.

\section{Conclusion}

Over five years, no-till with residues retained raised topsoil organic carbon
by \res{p4-dsoc-S1-NTR1}, \res{p4-dsoc-S2-NTR1} and
\res{p4-dsoc-S3-NTR1}~\si{\gram\per\kilo\gram} at S1, S2 and S3, while the
yield response to tillage depended on site and on residue management. Where
residues cannot be retained, no-till alone is unlikely to deliver either the
carbon or the yield benefit. All data are synthetic and illustrate the
analysis only.
